\documentclass{article}
%  ╔══════════════════════════════════════════════════════════╗
%  ║                         Title                            ║
%  ╚══════════════════════════════════════════════════════════╝
%NOTE: I want \title,\author to be near top, and this \renewcommand{\maketitle} fails if these are above it.
\usepackage[utf8]{inputenc}         % Used to compile any UTF-8 Character (and supports Unicode)
\usepackage{titling}

\renewcommand{\maketitle}{
    \quad
    \vspace{1em}
  \begin{center}
      \LARGE\bfseries \thetitle \par
      \vspace{0.6em}
      \large
  \end{center}
    \vspace{1em}
}

\title{Langlands for SL2}
\author{Nathanael Chwojko-Srawley}
\date{\today}

%  ╔══════════════════════════════════════════════════════════╗
%  ║                         packages                         ║
%  ╚══════════════════════════════════════════════════════════╝

\usepackage{extarrows}         % Used to compile any UTF-8 Character (and supports Unicode)
\usepackage{stackengine}
\usepackage{colortbl}
\usepackage{scalerel}
\usepackage{amsmath, amssymb} % Used to write better mathematical expressions (ex. \mathbb{R})
\usepackage{stackrel}         % for left-right arrows, staking symbol above and bellow
\usepackage{mathrsfs}
\usepackage{bbm}                     % for mathbbm{1}
\usepackage{euscript}
\usepackage{commutative-diagrams}
\usepackage[font=itshape]{quoting}
\usepackage{mathtools}
  \usepackage{enumitem}
\usepackage{amsthm}           % Used to create theorem enviroments.
\usepackage{graphicx}         % Let's you use images in LaTeX; ex the \includegraphics command.
\usepackage{hhline}
%\usepackage{luatexja}
%\usepackage{fontspec}
\usepackage{dsfont}
\usepackage{wasysym}          % emoji's!
\usepackage{float}            % package with ``H'' for figures and tables, ex. Images, tables, etc
%\usepackage{subcaption}      % More options with sub-captions.
\usepackage{setspace}         % More flexibility on spacing in document.
\usepackage{hyperref}         % let's me put hyperlinks
\usepackage{cleveref}         % let's me put hyperlinks
%\usepackage{tikz}             % for commutative diagrams https://tex.stackexchange.com/questions/419221/how-to-do-the-pushout-with-universal-property
\usepackage{tikz-cd}          % for commutative diagrams https://tex.stackexchange.com/questions/419221/how-to-do-the-pushout-with-universal-property
%\usepackage{quiver}           % for better commutative diagrams
\usepackage{bookmark}         % creating heading shortcut bar on the left as in word
\usepackage{longtable}        % create tables that extend past one page
%\usepackage{siunitx}         % required for alignment
%\usepackage{multirow}        % for multiple rows
\usepackage{microtype}        % makes nice text. Fixes some under-/over- filled box problems.
%\usepackage{fancyvrb}        % Let's you write code (different style than listings).
\usepackage[most]{tcolorbox} 	      % Makes nice boxes for thms, cor, prop, or anything else.
%\usepackage{enumitem}        % More flexibility on enumeration (don't know much).
\usepackage{xcolor}           % Colour text.
\usepackage{wrapfig}          % To wrap text around figure
\usepackage{listings}         % Create nice colour code blocks (customized in preamble)
\usepackage{thmtools}         % Customize theorem environment (in preamble).
%\usepackage{marvosym}        % Provide ``basic'' symbols (hazard, religious, smiley face etc.)
\usepackage{array}            % \begin{table}{>{\em}c c} -> 1st column italic (bfseries for bold)
%\usepackage[english]{babel}  % If I want to blindtext in english.
\usepackage{blindtext}
\usepackage[a4paper, total={6in, 8.5in}]{geometry}
%\usepackage[symbol]{footmisc} % For daggers and asterisk for footnotes. 1: *, 2: dagger, etc.
\usepackage{fancyhdr}
\usepackage{titlesec}
%\usepackage{pst-plot}          %To make nice plots: https://tex.stackexchange.com/questions/47594/how-can-i-make-a-graph-of-a-function#47596
% \usepackage[answerdelayed]{exercise}
\usepackage{etoolbox}          % for Dynkin diagrams
\usepackage{dynkin-diagrams}   % for Dynkin diagrams

% ╔═════════════════════════════════════════════════════╗
% ║           for figures via inkspace                  ║
% ╚═════════════════════════════════════════════════════╝

% to import figures via: https://github.com/gillescastel/inkscape-figures
\usepackage{import}
\usepackage{pdfpages}
\usepackage{transparent}


% This command is the ONLY command imported via the packages snippet, think of it as a tiny package written out
\newcommand{\incfig}[2][1]{%
    \def\svgwidth{#1\columnwidth}
    \import{./figures/}{#2.pdf_tex}
}

% ╔═══════════════════════════════════════════════════════════╗
% ║                 bibliography and citing                   ║
% ╚═══════════════════════════════════════════════════════════╝

\usepackage{csquotes}
% \usepackage[backend=biber,citestyle=alphabetic,bibstyle=authortitle]{biblatex}
\usepackage{biblatex}

\addbibresource{bibliography.bib}

% ╔═══════════════════════════════════════════════════════════╗
% ║                    colours (colors)                       ║
% ╚═══════════════════════════════════════════════════════════╝

\definecolor{dkgreen}{rgb}{0,0.6,0}
\definecolor{gray}{rgb}{0.5,0.5,0.5}
\definecolor{mauve}{rgb}{0.58,0,0.82}
\definecolor{lightyellow}{rgb}{1,1,0.6}
\definecolor{reddish}{HTML}{A01A3A}

%  ╔══════════════════════════════════════════════════════════╗
%  ║                       book format                        ║
%  ╚══════════════════════════════════════════════════════════╝

\pagestyle{fancy}

\setlength\headheight{20pt}

\renewcommand{\sectionmark}[1]{\markright{\thesection.\ #1}}

\fancyfoot[C,CO]{\textcolor{gray}{Nathanael Chwojko-Srawley}}
\fancyfoot[R,RO]{\thepage}{}

% Create a new page style for the first page
\fancypagestyle{firstpage}{
  \fancyhf{} % Clear header and footer
  \fancyhead[L]{\theauthor}
  \fancyhead[R]{\today} % Date in the top-right
  \fancyfoot[C]{\textcolor{gray}{Nathanael Chwojko-Srawley}}
  \fancyfoot[R]{\thepage}
}

% Apply the first page style conditionally
\AtBeginDocument{\thispagestyle{firstpage}}

%have format the title command
\newcommand{\chapter}[1]{%
  \clearpage
  \vspace*{30pt} % Space before the chapter title
  {\normalfont\huge\bfseries #1\par} % Chapter title formatting
  \vspace{20pt} % Space after the chapter title
}


\setlength{\parindent}{0pt} % don't like indentation infront of paragraphs
\setlength{\parskip}{1ex}   %so that there is still space between paragarphs


%  ╔══════════════════════════════════════════════════════════╗
%  ║                       Shortcuts                          ║
%  ╚══════════════════════════════════════════════════════════╝

%\ExerciseLevelInToc

% I like original mathcal better, but need lowercase.
\DeclareFontFamily{OT1}{pzc}{}
\DeclareFontShape{OT1}{pzc}{m}{it}{<-> s * [1.10] pzcmi7t}{}
\DeclareMathAlphabet{\mathpzc}{OT1}{pzc}{m}{it}


% mathbb
\newcommand{\A}{{\mathbb{A}}}
\newcommand{\C}{{\mathbb{C}}}
\newcommand{\F}{{\mathbb{F}}}
\newcommand{\N}{{\mathbb{N}}}
\newcommand{\Q}{{\mathbb{Q}}}
\newcommand{\R}{{\mathbb{R}}}
\newcommand{\V}{{\mathbb{V}}}
\newcommand{\Z}{{\mathbb{Z}}}
\newcommand{\BI}{{\mathbb{I}}}
\renewcommand{\H}{{\mathbb{H}}}
\renewcommand{\P}{{\mathbb{P}}}


% mathcal
% \newcommand{\co}{{\mathfrak{o}}} %mathcal{o} looks like wreath product, so I changed it to mathfrak
\newcommand{\co}{{\mathpzc{o}}}
\newcommand{\CA}{{\mathcal{A}}}
\newcommand{\CB}{{\mathcal{B}}}
\newcommand{\CC}{{\mathcal{C}}}
\newcommand{\CD}{{\mathcal{D}}}
\newcommand{\CF}{{\mathcal{F}}}
\newcommand{\CG}{{\mathcal{G}}}
\newcommand{\CH}{{\mathcal{H}}}
\newcommand{\CI}{{\mathcal{I}}}
\newcommand{\CK}{{\mathcal{K}}}
\newcommand{\CL}{{\mathcal{L}}}
\newcommand{\CM}{{\mathcal{M}}}
\newcommand{\CN}{{\mathcal{N}}}
\newcommand{\CO}{{\mathcal{O}}}
\newcommand{\CQ}{{\mathcal{Q}}}
\newcommand{\CR}{{\mathcal{R}}}
\newcommand{\CS}{{\mathcal{S}}}
\newcommand{\CT}{{\mathcal{T}}}
\newcommand{\CV}{{\mathcal{V}}}
\newcommand{\CW}{{\mathcal{W}}}
\newcommand{\CZ}{{\mathcal{Z}}}
% EXCPETION:
\newcommand{\CP}{{\mathcal{P}}}
% \newcommand{\CP}{\mathbb{C}\mathrm{P}}
\newcommand{\RP}{\mathbb{R}\mathrm{P}}


%Euscript
\newcommand{\EB}{\EuScript{B}}
\newcommand{\EC}{{\EuScript{C}}}
\newcommand{\EF}{{\EuScript{F}}}
\newcommand{\EL}{{\EuScript{L}}}
\newcommand{\EO}{{\EuScript{O}}}


%mathfrak (lowercase)
\newcommand{\fa}{{\mathfrak{a}}}
\newcommand{\fb}{{\mathfrak{b}}}
\newcommand{\fc}{{\mathfrak{c}}}
\newcommand{\fd}{{\mathfrak{d}}}
\newcommand{\fm}{{\mathfrak{m}}}
\newcommand{\fn}{{\mathfrak{n}}}
\newcommand{\fo}{{\mathfrak{o}}}
\newcommand{\fp}{{\mathfrak{p}}}
\newcommand{\fq}{{\mathfrak{q}}}


%mathfrak (upper case)
\newcommand{\FA}{{\mathfrak{A}}}
\newcommand{\FB}{{\mathfrak{B}}}
\newcommand{\FC}{{\mathfrak{C}}}
\newcommand{\FD}{{\mathfrak{D}}}
\newcommand{\FK}{{\mathfrak{K}}}
\newcommand{\FM}{{\mathfrak{M}}}
\newcommand{\FN}{{\mathfrak{N}}}
\newcommand{\FO}{{\mathfrak{O}}}
\newcommand{\FP}{{\mathfrak{P}}}


%mathscr
\newcommand{\ff}{{\mathscr{F}}}
\newcommand{\SA}{\mathscr{A}}
\newcommand{\SC}{\mathscr{C}}
\newcommand{\SF}{\mathscr{F}}
\newcommand{\SG}{\mathscr{G}}
\newcommand{\SH}{\mathscr{H}}
\newcommand{\SI}{\mathscr{I}}
\newcommand{\SK}{\mathscr{K}}
\newcommand{\SN}{\mathscr{N}}
\newcommand{\SQ}{\mathscr{Q}}
\newcommand{\SR}{\mathscr{R}}
\newcommand{\SV}{\mathscr{V}}
\newcommand{\SZ}{\mathscr{Z}}


% operatorname -- 2 letters (lowercase and upper case)
\newcommand{\ab}{{\operatorname{ab}}}
\newcommand{\ad}{{\operatorname{ad}}}
\newcommand{\ch}{{\operatorname{ch}}}
\newcommand{\ev}{{\operatorname{ev}}}
\newcommand{\id}{{\operatorname{id}}}
\newcommand{\im}{{\operatorname{im}}}
\newcommand{\nm}{{\operatorname{Nm}}}
\newcommand{\op}{{\operatorname{op}}}
\newcommand{\rk}{{\operatorname{rk}}}
\newcommand{\sh}{{\operatorname{sh}}}
\newcommand{\tr}{{\operatorname{tr}}}

\newcommand{\Ab}{{\operatorname{Ab}}}
\newcommand{\Br}{{\operatorname{Br}}}
\newcommand{\Ch}{{\operatorname{Ch}}}
\newcommand{\Cl}{{\operatorname{Cl}}}
\newcommand{\GL}{{\operatorname{GL}}}
\newcommand{\Nm}{{\operatorname{Nm}}}
\newcommand{\SL}{{\operatorname{SL}}}

\renewcommand{\Re}{{\operatorname{Re}}}
\renewcommand{\Im}{{\operatorname{Im}}}


%categories
\newcommand{\Grp}{{\textbf{Grp}}}
\newcommand{\Fld}{{\textbf{Fld}}}
\newcommand{\Sh}{{\textbf{Sh}}}
\newcommand{\Set}{{\textbf{Set}}}
\newcommand{\Mod}{{\textbf{Mod}}}
\newcommand{\PSh}{{\textbf{PSh}}}
\newcommand{\Sch}{{\textbf{Sch}}}
\newcommand{\Top}{{\textbf{Top}}}
\newcommand{\RgSp}{{\textbf{RgSp}}}
\newcommand{\Ring}{{\textbf{Ring}}}


%operatorname -- 3+ letters (lowercase and upper case)
\newcommand{\rad}{{\operatorname{rad}}}
\newcommand{\sep}{{\operatorname{sep}}}
\newcommand{\res}{{\operatorname{res}}}
\newcommand{\sgn}{{\operatorname{sgn}}}
\newcommand{\pre}{{\operatorname{pre}}}
\newcommand{\ord}{{\operatorname{ord}}}
\newcommand{\vol}{{\operatorname{vol}}}
\newcommand{\lcm}{{\operatorname{lcm}}}
\newcommand{\len}{{\operatorname{len}}}

\newcommand{\conj}{{\operatorname{conj}}}
\newcommand{\disc}{{\operatorname{disc}}}
\newcommand{\stab}{{\operatorname{stab}}}
\newcommand{\kdim}{{\operatorname{kdim}}}
\newcommand{\diag}{{\operatorname{diag}}}

\newcommand{\trdeg}{\operatorname{trdeg}}
\newcommand{\coker}{{\operatorname{coker}}}
\newcommand{\codim}{{\operatorname{codim}}}

\newcommand{\Ann}{{\operatorname{Ann}}}
\newcommand{\Aut}{{\operatorname{Aut}}}
\newcommand{\Der}{{\operatorname{Der}}}
\newcommand{\Div}{{\operatorname{Div}}}
\newcommand{\Emb}{{\operatorname{Emb}}}
\newcommand{\End}{{\operatorname{End}}}
\newcommand{\Ext}{{\operatorname{Ext}}}
\newcommand{\Gal}{{\operatorname{Gal}}}
\newcommand{\Hom}{{\operatorname{Hom}}}
\newcommand{\Ind}{{\operatorname{Ind}}}
\newcommand{\Inn}{{\operatorname{Inn}}}
\newcommand{\Int}{{\operatorname{int}}}
\newcommand{\Irr}{{\operatorname{Irr}}}
\newcommand{\Jac}{{\operatorname{Jac}}}
\newcommand{\Mor}[1]{\operatorname{Mor}(\textbf{#1})}
\newcommand{\Nil}{{\operatorname{Nil}}}
\newcommand{\Obj}{{\operatorname{Obj}}}
\newcommand{\Out}{{\operatorname{Out}}}
\newcommand{\Pic}{{\operatorname{Pic}\,}}
\newcommand{\Rep}{{\operatorname{Rep}}}
\newcommand{\Res}{{\operatorname{Res}}}
\newcommand{\Spl}{{\operatorname{Spl}}}
\newcommand{\Syl}{{\operatorname{Syl}}}
\newcommand{\Sym}{{\operatorname{Sym}}}
\newcommand{\Tor}{{\operatorname{Tor}}}

\newcommand{\Char}{{\operatorname{char}}}
\newcommand{\Frac}{{\operatorname{Frac}}}
\newcommand{\Frob}{{\operatorname{Frob}}}
\newcommand{\Proj}{{\operatorname{Proj}\,}}
\newcommand{\RMod}{{}_R\operatorname{Mod}}
\newcommand{\SExt}{{\mathcal{E}\emph{xt}}}
\newcommand{\SHom}{{\emph{Hom}}}
\newcommand{\Span}{{\operatorname{span}}}
\newcommand{\Spec}{{\operatorname{Spec}\,}}
\newcommand{\Supp}{{\operatorname{Supp}}}
\newcommand{\Tens}[1]{#1\otimes --}
\newcommand{\fHom}[1]{\operatorname{Hom}(#1, --)}

\newcommand{\mSpec}{{\operatorname{mSpec}}}
\newcommand{\cofHom}[1]{\operatorname{Hom}(--, #1)}


% connectors
% \renewcommand{\mid}{\big|}
\newcommand{\sse}{\subseteq}
\newcommand{\subg}{\leqslant}
\newcommand{\supg}{\geqslant}
\newcommand{\ssne}{\subsetneq}
\newcommand{\isosubg}{\lesssim}
\newcommand{\nsse}{\not\subseteq}
\newcommand{\nsubg}{\trianglelefteq}
\newcommand{\nsupg}{\trianglerighteq}
\newcommand{\csubg}{\blacktriangleleft}
\renewcommand{\mid}{\big|}


% Arrows
\newcommand{\rw}{\rightarrow}
\newcommand{\Rw}{\Rightarrow}
\newcommand{\lw}{\leftarrow}
\newcommand{\Lw}{\Leftarrow}
\newcommand{\lrw}{\leftrightarrow}
\newcommand{\LRw}{\Leftrightarrow}
\newcommand{\acton}{\curvearrowright}
\newcommand{\actson}{\curvearrowright}
\newcommand\mapsfrom{\mathrel{\reflectbox{\ensuremath{\mapsto}}}}

% arrows for commutative diagram: create four new angled double-struck arrows
\newcommand\myrot[1]{\mathrel{\rotatebox[origin=c]{#1}{$\Rightarrow$}}}
\newcommand\NEarrow{\myrot{45}}
\newcommand\NWarrow{\myrot{135}}
\newcommand\SWarrow{\myrot{-135}}
\newcommand\SEarrow{\myrot{-45}}


%shorthands
\newcommand{\oln}[1]{{\overline{#1}}}
\renewcommand{\bf}[1]{\textbf{#1}}
\newcommand{\qn}{\quad\newline}


% limits
\newcommand{\Lim}{{\varprojlim}}
\newcommand{\coLim}{{\varinjlim}}
\newcommand{\Dlim}{\lim\limits_{\longrightarrow}}
\newcommand{\coDlim}{\lim\limits_{\longleftarrow}}


%for saturated ideals in the localization section (I don't like these, I think I'll delete)
\newcommand{\LIm}{{}^e }
\newcommand{\LPre}{{}^c }

%  ╔══════════════════════════════════════════════════════════╗
%  ║                    Custom Commands                       ║
%  ╚══════════════════════════════════════════════════════════╝

% swap varphi and phi
\let\temp\phi
\let\phi\varphi
\let\varphi\temp

\newcommand{\placeholder}{\_\_\_\_\_}

\newcommand{\set}[2]{\left\{#1 \ \middle|\ #2\right\}}

%mod should not have the crazy amount of space to its right!
\renewcommand{\mod}{\,\operatorname{mod}\,}

% dcups
\newcommand{\dcup}{\sqcup}
\newcommand{\Dcup}{\coprod}
% \newcommand{\Dcup}{\bigsqcup}

\newcommand{\nbot}{\mathrel{\rlap{$\bot$}\mkern2mu/}}

%a nice large circle
\newcommand{\tikzcircle}[2][red,fill=black]{\tikz[baseline=-0.5ex]\draw[#1,radius=#2] (0,0) circle ;}%

\makeatletter
\newcommand*\bigcdot{\mathpalette\bigcdot@{.5}}
\newcommand*\bigcdot@[2]{\mathbin{\vcenter{\hbox{\scalebox{#2}{$\m@th#1\bullet$}}}}}
\makeatother

%somehow not a default command
\def\rddots{\mathstrut^{.^{.^{.}}}}

%% new delimiter
\DeclarePairedDelimiter{\ceil}{\lceil}{\rceil}


%  ╔══════════════════════════════════════════════════════════╗
%  ║                     environments                         ║
%  ╚══════════════════════════════════════════════════════════╝


%remember the option: breakable

% create tcolor boxes
\newtcbtheorem[number within=section]{thm}{Theorem}%
{
colback=green!5,
colframe=green!35!black,
fonttitle=\bfseries,
label type=theorem
}{th}

\newtcbtheorem[number within=section, use counter from=thm]{lem}{Lemma}%
{
colback=blue!5,
colframe=black!35!black,
fonttitle=\bfseries,
label type=lemma
}{lm}
\newtcbtheorem[number within=section, use counter from=thm]{prop}{Proposition}%
{
colback=white!5,
colframe=white!35!black,
fonttitle=\bfseries,
label type=proposition
}{pr}
\newtcbtheorem[number within=section, use counter from=thm]{cor}{Corollary}%
{colback=blue!5,
colframe=blue!35!black,
fonttitle=\bfseries,
label type=corollary
}{co}
\newtcbtheorem[number within=section, use counter from=thm]{axiom}{Axiom}%
{colback=black!5,
colframe=yellow!35!black,
fonttitle=\bfseries,
label type=axiom
}{ax}
\newtcbtheorem[number within=section, use counter from=thm]{defn}{Definition}%
{colback=red!5,
colframe=red!35!black,
fonttitle=\bfseries,
label type=definition
}{df}


% for <s-cr> to create \item[$\LRw$] instead of \item
\newenvironment{equivEnumerate}
 {\begin{enumerate}
	}{
\end{enumerate}}



%Custom enviroments
 \newenvironment{note}
 {\begin{tcolorbox}[enhanced, sharp corners, colback=white , borderline={0.2pt}{0pt}{black}]
   \begin{quote}\textbf{Note}:
   }{
   \end{quote}\end{tcolorbox}}

 \newenvironment{intuition}
 {\begin{quote}\textbf{Intuition}:
   }{
   \end{quote}}

\newtcbtheorem[number within=section, use counter from=thm]{titledBox}{}{%
  enhanced,
  sharp corners,
  colback=white,
  colframe=black,
  borderline={0.2pt}{0pt}{black},
  left=2mm,
  right=2mm,
  top=2mm,
  bottom=2mm,
  title style={opacity=1},
  attach title to upper,
  before upper={\textbf{\tcbtitletext}\quad},
  coltitle=black,
  fonttitle=\bfseries,
  separator sign={\quad}
}{box}


%better example and proof environments
\def\exampletext{Example} % If English
\colorlet{colexam}{red!55!black} % Global example color


\newtcbtheorem[number within=section, use counter from=thm]{example}{Example}{% \newtcolorbox[use counter=testexample]{testexamplebox}{%
% Example Frame Start
empty,% Empty previously set parameters
title={\exampletext \thetcbcounter #1},% use \thetcbcounter to access the testexample counter text
% Attaching a box requires an overlay
attach boxed title to top left,
   % Ensures proper line breaking in longer titles
   minipage boxed title,
% (boxed title style requires an overlay)
boxed title style={empty,size=minimal,toprule=0pt,top=4pt,left=3mm,overlay={}},
coltitle=colexam,fonttitle=\bfseries,
before=\par\medskip\noindent,parbox=false,boxsep=0pt,left=3mm,right=0mm,top=2pt,breakable,pad at break=0mm,
   before upper=\csname @totalleftmargin\endcsname0pt, % Use instead of parbox=true. This ensures parskip is inherited by box.
% Handles box when it exists on one page only
overlay unbroken={\draw[colexam,line width=.5pt] ([xshift=-0pt]title.north west) -- ([xshift=-0pt]frame.south west); },
% Handles multipage box: first page
overlay first={\draw[colexam,line width=.5pt] ([xshift=-0pt]title.north west) -- ([xshift=-0pt]frame.south west); },
% Handles multipage box: middle page
overlay middle={\draw[colexam,line width=.5pt] ([xshift=-0pt]frame.north west) -- ([xshift=-0pt]frame.south west); },
% Handles multipage box: last page
overlay last={\draw[colexam,line width=.5pt] ([xshift=-0pt]frame.north west) -- ([xshift=-0pt]frame.south west); }%
}{ex}



\def\prooftext{Proof} % If English
\NewDocumentEnvironment{Proof}{ O{} }
{
    \colorlet{colexam}{black!55!black} % Global example color
    \newtcolorbox[number within=section]{testproofbox}{%
        empty,% Empty previously set parameters
        title={\emph{\prooftext} \thetcbcounter: #1},
        attach boxed title to top left,
        minipage boxed title,
        boxed title style={empty,size=minimal,toprule=0pt,top=4pt,left=3mm,overlay={}},
        coltitle=colexam,
        fonttitle=\bfseries,
        before=\par\medskip\noindent,
        parbox=false,
        boxsep=0pt,
        left=3mm,
        right=0mm,
        top=2pt,
        breakable,
        pad at break=0mm,
        before upper=\csname @totalleftmargin\endcsname0pt,
        % Removed height constraints
        % Added flexible width
        width=\dimexpr\textwidth-3mm\relax,
        % Modified overlays
        overlay unbroken={\draw[colexam,line width=.5pt] ([xshift=-0pt]title.north west) -- ([xshift=-0pt]frame.south west); },
        overlay first={\draw[colexam,line width=.5pt] ([xshift=-0pt]title.north west) -- ([xshift=-0pt]frame.south west); },
        overlay middle={\draw[colexam,line width=.5pt] ([xshift=-0pt]frame.north west) -- ([xshift=-0pt]frame.south west); },
        overlay last={\draw[colexam,line width=.5pt] ([xshift=-0pt]frame.north west) -- ([xshift=-0pt]frame.south west); },%
    }
    \begin{testproofbox}
}
{\end{testproofbox}\endlist}


%  ╔══════════════════════════════════════════════════════════╗
%  ║                    for Dynkin diagrams                   ║
%  ╚══════════════════════════════════════════════════════════╝

\def\row#1/#2!{#1_{\IfStrEq{#2}{}{n}{#2}} & \dynkin{#1}{#2}\\}
\newcommand{\tble}[1]{
   \renewcommand*\do[1]{\row##1!}
   \[
      \begin{array}{ll}\docsvlist{#1}\end{array}
   \]
}

%  ╔══════════════════════════════════════════════════════════╗
%  ║                         links                            ║
%  ╚══════════════════════════════════════════════════════════╝

\definecolor{LinkColour}{HTML}{A64A3B} % favourite
% \definecolor{LinkColour}{HTML}{8C3D32} % 2nd place
\hypersetup{
  colorlinks,
  citecolor=black,
  filecolor=magenta,
  linkcolor=LinkColour,
  urlcolor=blue,
}


%  ╔══════════════════════════════════════════════════════════╗
%  ║                          misc                            ║
%  ╚══════════════════════════════════════════════════════════╝

\stackMath %to be able to stack text in math mode
\makeindex

%import options: all, bibliography, book-formatting, booksSetting, customEnvironments, dynkin, hw-formatting, language-setup, newcommands, packages, preamble, proofs, settings, tcolorbox, test.svg,
\begin{document}
\maketitle

\section{Abstract}

Let $F$ be a locally compact nonarchimedean field with residue characteristic $p$ (e.g. finite extensions of
$\mathbb{Q}_p$), and let $K$ be an algebraically closed field of characteristic $\ne p$. This paper analyzes
smooth $K$-representations of $\SL_2(F)$ via restriction from $\GL_2(F)$. The reader can assume $K
= \mathbb{C}$ to link this to the intuition of complex representations like
in~\cite{bushnellLocalLanglandsConjecture2006}. While the Local Langlands correspondence for $\GL_n$ is
a bijection, the restriction to $\SL_2$ creates fibers of size $1$, $2$, or $4$. This paper provides a structural proof of the size of these fibers—known as \emph{L-packets} or \emph{L-indistinguishable classes}—using Clifford Theory adapted for $p$-adic groups and the uniqueness of Whittaker models.

\begin{titledBox}{Inspiration}{insp}
  This paper is completely inspired and modeled by~\cite{henniartRepresentationsSL2F2025}.
  This is my own re-writing of the first sections of Henniart's and Vign\'eras paper with details useful to
  me (and a few math typos fixed).
\end{titledBox}

\section{Build up}

Let us start with some useful results from representation theory:

\begin{thm}{Clifford's Theorem}{cliffordThm}
	Let $G$ be a group, $N \trianglelefteq G$ a normal subgroup, $K$ a field, and $\pi: G \to \GL(V)$ an irreducible
	representation. Then the restriction $\pi|_N$ is semisimple and breaks up into a direct sum of irreducible representations
	of $N$ of equal dimensions. These irreducible representations of $N$ lie in a single orbit under the action of $G$ by conjugation.
\end{thm}

\begin{Proof}
	Let $W \subseteq V$ be an irreducible subrepresentation of $\pi|_N$ (assuming $V$ is finite-dimensional, such a $W$ exists).
	For any $g \in G$, define the subspace $gW = \{ \pi(g)w \mid w \in W \}$. I claim that $gW$ is an irreducible
	subrepresentation of $N$.

	For any $n \in N$ and $v' \in gW$ (where $v' = \pi(g)w$), the action is:
	\[ \pi(n)v' = \pi(n)\pi(g)w = \pi(g) \pi(g^{-1}ng) w \]
	Since $N$ is normal, $n' = g^{-1}ng \in N$. Since $W$ is an $N$-submodule, $\pi(n')w \in W$.
	Thus $\pi(n)v' \in gW$, so $gW$ is $N$-invariant.

	Furthermore, the map $\pi(g): W \to gW$ is a vector space isomorphism. Thus, $\dim(gW) = \dim(W)$.
	Since $W$ is irreducible, $gW$ must also be irreducible (any proper submodule of $gW$ would pull back to a proper submodule of $W$).

	Consider the sum of these subspaces $V' = \sum_{g \in G} gW$.
	For any $h \in G$, $\pi(h)V' = \sum_{g \in G} \pi(h)gW = \sum_{g \in G} (hg)W$.
	Since $hg$ runs over $G$ as $g$ does, $V'$ is $G$-invariant.
	Since $V$ is irreducible and $V' \neq 0$, we must have $V' = V$.

	Therefore, $V$ is a sum of irreducible $N$-subrepresentations ($gW$), implying $\pi|_N$ is semisimple.
	By construction, every irreducible component is of the form $gW$, which implies they are all conjugate
	(lie in one orbit) and have dimension equal to $\dim(W)$.
\end{Proof}

Thus, assume $H\nsubg G$ is a finite index normal subgroup, and let $\pi$ be an irreducible $K$-representation
of $H$. Let us show that it is a component of a $K$-representation $\Pi$ of $G$ whose restriction to $H$
contains $\pi$. First, induct so that we get $\Ind_H^G\pi$. This is finite length still as $[G:H] < \infty$,
and contains an irreducible representation, say $\Pi$. Then there is a non-trivial map $\Pi \hookrightarrow \Ind_H^G
\pi$. By Frobenius:
\[
  \Hom_G(\Pi, \Ind_H^G \pi) \cong \Hom_H(\Pi|_H, \pi)
\]
Hence, as $\pi$ is irreducible and the restriction is semisimple, $\pi$ must be in the direct summands of
$\Pi|_H$. Next, if $\chi$ is a $K$-character of $G$ that is trivial on $H$, the restriction of $\Pi\otimes
\chi$ to $H$ contains $\pi$, for
\[
  (\Pi \otimes \chi)|_H \cong \Pi|_H
\]
If there is another irreducible $K$-representation $\Pi'$ that contains $\pi$, it is a twist of $\Pi$:

\begin{lem}{Gallagher's Lemma}{lem1}
    Let $K$ be an algebraically closed field, $G$ a group, and $H \trianglelefteq G$ a normal subgroup such that the quotient $A = G/H$ is a finite abelian group.
    Let $\Pi$ be an irreducible representation of $G$. If $\Pi'$ is an irreducible representation of $G$ such that
    \[ \mathrm{Hom}_H(\Pi'|_H, \Pi|_H) \neq 0 \]
    (i.e. their restrictions overlap), then $\Pi' \simeq \Pi \otimes \chi$ for some character $\chi$ of $A$ (lifted to $G$).
\end{lem}

\begin{Proof}
    By assumption, $\mathrm{Hom}_H(\Pi'|_H, \Pi|_H) \neq 0$. Since $G/H$ is finite, Induction is both the left and right adjoint to Restriction. Using Frobenius Reciprocity:
    \[ 0 \neq \mathrm{Hom}_H(\Pi'|_H, \Pi|_H) \cong \mathrm{Hom}_G(\Pi', \mathrm{Ind}_H^G(\Pi|_H)) \]
    Thus, $\Pi'$ is isomorphic to a subrepresentation of $\mathrm{Ind}_H^G(\Pi|_H)$.

    By the Tensor Identity (Projection Formula), we have:
    \[ \mathrm{Ind}_H^G(\Pi|_H) \cong \mathrm{Ind}_H^G(1_H \otimes \Pi|_H) \cong (\mathrm{Ind}_H^G 1_H) \otimes \Pi \]
    Since $G/H$ is a finite abelian group and $K$ is algebraically closed, the regular representation of the quotient decomposes into 1-dimensional characters:
    \[ \mathrm{Ind}_H^G 1_H \cong \bigoplus_{\chi \in \widehat{G/H}} \chi \]
    Substituting this back, we get:
    \[ \mathrm{Ind}_H^G(\Pi|_H) \cong \left( \bigoplus_{\chi} \chi \right) \otimes \Pi \cong \bigoplus_{\chi} (\chi \otimes \Pi) \]
    Since $\Pi'$ is irreducible and is a subrepresentation of this direct sum, it must be isomorphic to one of the summands. Note that since $\Pi$ is irreducible and $\chi$ is 1-dimensional, $\chi \otimes \Pi$ is also irreducible. Therefore, $\Pi' \simeq \chi \otimes \Pi$ for some character $\chi$.
\end{Proof}

\begin{prop}{Clifford Theory for Smooth Representations}{basicProp}
  Let $H\nsubg G$ be normal with a compact quotient $G/H$ and assume that $V$ is smooth (so $V|_H$ is smooth).

  \begin{enumerate}
    \item  If $V$ is \emph{finitely generated} then $V|_H$ is finitely generated %[Henniart 2001, lemme 4.1].

  \item If $V$ is \emph{irreducible}, any irreducible subrepresentation of $V|_H$ (when there exists one)
    extends to a (smooth and irreducible) representation of an open subgroup of $G$ of finite index which is
    admissible if $V$ is as well% [loc. cit., proposition 4.4].

  \item If $V$ is \emph{irreducible} and $V|_H$ contains an irreducible subrepresentation or is
    \emph{noetherian} (any subrepresentation is finitely generated), then $V|_H$ is semisimple of finite
    length %[loc. cit., théorème 4.2].
  \end{enumerate}
\end{prop}

\begin{Proof}
    \begin{enumerate}
      \item
        Let $\{v_1, \dots, v_n\}$ be a generating set for $V$ as a $G$-module. Since $G/H$ is compact, there exists a
        compact set $C \subseteq G$ such that $G = H \cdot C$ (i.e., $C$ surjects onto the quotient).
        Any element $v \in V$ can be written as a finite sum $\sum g_i v_{j_i}$. Writing $g_i = h_i c_i$ with $h_i \in H, c_i \in C$,
        we see that $V$ is generated as an $H$-module by the subspace $W = \mathrm{span}_{\mathbb{C}} \{ c \cdot v_k \mid c \in C, 1 \le k \le n \}$.
        Since $V$ is smooth, the stabilizer of each $v_k$ is open. Since $C$ is compact, the image $C \cdot v_k$ lies in a finite dimensional
        subspace (smooth action of a compact set on a vector generates a finite dim space). Thus $W$ is finite dimensional.
        A finite generating set for $W$ generates $V$ over $H$.

      \item
        Let $W \subset V|_H$ be an irreducible subrepresentation. Let $S = \mathrm{Stab}_G(W) = \{g \in G \mid gW \subseteq W\}$.
        Since $W$ is irreducible and $H$ is normal, $S$ is the stabilizer of the equivalence class of $W$.
        By standard Clifford theory arguments (\textcolor{red}{cite here}) (adapted for smooth representations), $V$ is induced from a representation of $S$.
        Since $V$ is smooth, stabilizers are open, so $S$ is an open subgroup.
        If $G/H$ is compact, $G/S$ is finite (since $S$ contains $H$ and is open in the compact topology of $G/H$).

      \item
        This relies on the fact that representations of compact groups are semisimple (complete reducibility).
        Consider the action of the quotient group $K = G/H$. $V$ can be viewed somewhat like a representation of $K$.
        Since $K$ is compact, the unitary representations decompose into a direct sum of finite-dimensional irreducibles.
        Because $V|_H$ has an irreducible subrepresentation (or is Noetherian), we can form the sum of all conjugates $gW$.
        Since $G/H$ is compact, there are only finitely many non-isomorphic conjugates (Clifford theory implies the orbit is finite if the index is finite; compactness mirrors this).
        Thus $V|_H$ is a finite direct sum of isotypic components.
    \end{enumerate}
\end{Proof}

\begin{prop}{Extension of Admissible Representations}{admissibleExt}
   Let $W$ be an admissible irreducible $R$-representation of $H$. Assume:
  \begin{enumerate}[label=(A\arabic*)]
      \item Any finitely generated admissible $R$-representation of $G$ has finite length.
      \item Any finitely generated smooth $R$-representation of $H$ is noetherian.
  \end{enumerate}

  Then:
  \begin{enumerate}[label=(\arabic*)]
      \item If (A1) and (A2) are true, $W$ is contained in some irreducible admissible $R$-representation of $G$ restricted to $H$.
      \item If (A1) is true, $W$ is a quotient of some irreducible admissible $R$-representation of $G$ restricted to $H$.
  \end{enumerate}
\end{prop}

\begin{Proof}
  We prove (2) first, as (1) follows from it via duality.

  \textbf{Proof of (2):}
  Let $\mathcal{I}(W) = \operatorname{Ind}_H^G W$. Since $G/H$ is compact and $W$ is admissible and finitely generated (being irreducible), $\mathcal{I}(W)$ is admissible and finitely generated.
  By assumption (A1), $\mathcal{I}(W)$ has finite length. Therefore, it must contain an irreducible subrepresentation $V$.
  \[ 0 \neq \operatorname{Hom}_G(V, \operatorname{Ind}_H^G W) \]
  By Frobenius Reciprocity (Restriction is left adjoint to Induction), we have:
  \[ \operatorname{Hom}_G(V, \operatorname{Ind}_H^G W) \cong \operatorname{Hom}_H(V|_H, W) \]
  Thus, there exists a non-zero map $V|_H \to W$. Since $W$ is irreducible, this map is surjective. Hence, $W$ is a quotient of $V|_H$.

  \textbf{Proof of (1):}
  Let $\check{W}$ be the smooth contragredient (dual) of $W$. Since $W$ is admissible, $\check{W}$ is also admissible and irreducible.
  We apply result (2) to $\check{W}$. There exists an irreducible admissible representation $V'$ of $G$ such that $\check{W}$ is a quotient of $V'|_H$:
  \[ V'|_H \twoheadrightarrow \check{W} \]
  Taking the smooth dual of this sequence reverses the arrows:
  \[ (\check{W})^\vee \hookrightarrow (V'|_H)^\vee \]
  Since $W$ is admissible, it is reflexive, so $(\check{W})^\vee \cong W$.
  Using the compatibility of duality and restriction (which relies on $H$ being open/normal and the Noetherian properties from (A2) to ensure well-behaved admissible duals over $R$), we have $(V'|_H)^\vee \cong (\check{V'})|_H$.
  Let $V = \check{V'}$. Then $W \hookrightarrow V|_H$, completing the proof.
\end{Proof}

\subsection*{Setup for Abelian Clifford Theory}

Suppose $H$ is a closed \emph{normal} subgroup of a \emph{locally profinite} group $G$.
Recall that if $V$ is a smooth representation of $G$, then $V|_H$ is smooth.
Furthermore, if $V|_H$ has finite length (as guaranteed by~\cref{pr:admissibleExt} when $V$ is admissible and $G/H$ is compact), we define the following structural components:

\begin{defn}{Clifford Components}{cliffordDefs}
Let $V$ be an irreducible representation of $G$ such that $V|_H$ has finite length.
\begin{itemize}
    \item Let $W \subseteq V|_H$ be an irreducible subrepresentation and $\pi$ its isomorphism class.
    \item Let $G_\pi = \{ g \in G \mid \pi^g \cong \pi \}$ be the \emph{stabilizer} of $\pi$ in $G$, where
      $\pi^g(h) = \pi(g^{-1}hg)$.
    \item Let $V_\pi$ be the \emph{$\pi$-isotypic component} of $V|_H$.
\end{itemize}
\end{defn}

Since $V$ is smooth, $\operatorname{Stab}_G(W)$ is open, implies that $G_\pi$ is an open subgroup of $G$.
Consequently, we have the fundamental Clifford decomposition:
\[
  V \cong \operatorname{Ind}_{G_\pi}^G(V_\pi)
\]
and the length formula:
\[
  \operatorname{length}(V|_H) = [G:G_\pi] \cdot \operatorname{length}(V_\pi|_H)
\]
\begin{lem}{Clifford Theory for Abelian Quotients}{lem2}
  Let $V$ be an irreducible representation of $G$, $H$ a closed subgroup, and assume $V|_H$ has finite length.
  Let $W$ be an irreducible subrepresentation of $V|_H$ and $\pi$ its isomorphism class.
  Let $X_V = \{ \chi \in \widehat{G/H} \mid V \otimes \chi \cong V \}$ be the group of self-twists.

  Assume that the quotient \(A = G/H\) is \emph{abelian}. Then:

  \begin{enumerate}
    \item \(G_{\pi}\) is a normal subgroup of \(G\), $G_\pi \nsubg G$ and does not depend on the choice of the constituent \(\pi\) in \(V|_{H}\).
    The smooth \(K\)-characters of \(G\) trivial on \(G_{\pi}\) are exactly the elements of \(X_V\).

    \item Assume \(K\) is algebraically closed.
    \begin{enumerate}
      \item Any irreducible subquotient of the smooth induction \(\mathrm{Ind}_H^G \mathbf{1}\) is a smooth character \(\chi\) of \(G\) trivial on \(H\).

      \item Any irreducible \(K\)-representation of \(G\) containing \(\pi\) (upon restriction to $H$) is isomorphic to \(V \otimes \chi\) for some smooth character \(\chi\) of \(G\) trivial on \(H\).

      \item \emph{Multiplicity 1:} If \(V|_H\) has multiplicity 1, then:
      \begin{itemize}
          \item \(W = V_{\pi}\) (the isotypic component coincides with the irreducible subrepresentation).
          \item For a smooth character \(\chi\) of \(G\) trivial on \(H\), \(V \otimes \chi \simeq V\) if and only if \(\chi\) is trivial on \(G_{\pi}\).
          \item \(G_{\pi}\) is the largest subgroup \(I\) (with $H \subseteq I \subseteq G$) such that \(\mathrm{length}(V|_I) = \mathrm{length}(V|_H)\).
      \end{itemize}

      \item \emph{Degree of Self-Twists:} If \(V|_H\) has multiplicity 1, the number of self-twists is:
      \[
        |X_V| = \begin{cases}
            [G:G_{\pi}] & \text{if } \mathrm{char}(K) = 0, \\
            [G:G_{\pi,\ell}] & \text{if } \mathrm{char}(K) = \ell > 0,
        \end{cases}
      \]
      where \(G_{\pi,\ell}\) is the smallest subgroup of \(G\) containing \(G_{\pi}\) such that the index \([G:G_{\pi,\ell}]\) is relatively prime to \(\ell\).
    \end{enumerate}
  \end{enumerate}
\end{lem}


\begin{Proof}
\begin{enumerate}
\item The isotypic components of \(\Pi|_H\) are \(G\)-conjugate, their \(G\)-stabilizers are \(G\)-conjugate and contain \(H\) hence they are equal because \(G/H\) is abelian. Since \(V \otimes \chi \simeq \mathrm{Ind}_{G_{\pi}}^G(\chi|_{G_{\pi}} \otimes V_{\pi})\) for any smooth \(K\)-character \(\chi\) of \(G\), the smooth \(K\)-characters of \(G\) trivial on \(G(\pi)\) are in \(X_V\).

\item
\begin{enumerate}
\item[(a)] For any closed subgroup \(Q\) of \(G\) and a smooth \(K\)-representation \(X\) of \(Q\), the representation \(\mathrm{Ind}_Q^G X\) is the space of functions \(f: G \to X\) with the property \(f(qgk) = qf(g)\) for \(q \in Q\), \(g \in G\), \(k \in K_f\) for some open subgroup \(K_f\) of \(G\), with the action of \(G\) by right translation, and where \(\mathrm{ind}_Q^G 1\) is the subrepresentation on the subspace of functions of compact support modulo \(Q\). When \(G/Q\) is compact, \(\mathrm{Ind}_Q^G X = \mathrm{ind}_Q^G X\).

Let \(V \supset U\) be \(G\)-stable subspaces with \(V/U\) irreducible. Suppose \(V\) generated by an element
\(f\) (indeed \(V'/U' \simeq V/U\) for the \(G\)-stable space \(V'\) generated by \(f \in V \setminus U\) and
the kernel \(U'\) of the map \(V' \to V/U\)). There is an open subgroup \(K\) of \(G\) which fixes \(f\). Then \(U \subset V \subset \mathrm{ind}_K^G 1\) and one is reduced to the case where \(G/H\) is finite.

\item[(b)] The proof of Lemma \ref{lm:lem1} remains valid with the smooth induction \(\mathrm{Ind}_H^G\), which is the smooth compact induction \(\mathrm{ind}_H^G 1\), because \(G/H\) is compact, so that \(\mathrm{ind}_H^G(\Pi|_H) = \Pi \otimes \mathrm{ind}_H^G 1\).

\item[(c)] Any smooth character \(\chi\) of \(G\) trivial on \(H\) with \(\mathrm{ind}_{G_{\pi}}^G(V_{\pi}) \simeq \mathrm{ind}_{G_{\pi}}^G(V_{\pi} \otimes \chi|_{G_{\pi}})\) is trivial on \(G_{\pi}\). Indeed, restricting to \(G_{\pi}\) we see that \(V_{\pi} \otimes \chi|_{G_{\pi}}\) is conjugate to \(V_{\pi}\) by some \(g \in G\). Restricting to \(H\) gives that \(\pi \simeq \pi^g\), so \(g \in G_{\pi}\), hence \(V_{\pi} \otimes \chi|_{G_{\pi}} \simeq V_{\pi}\). As \(\mathrm{Ker}(\chi)\) is open in \(G\) and \(G/H\) is compact, \(J = \mathrm{Ker}(\chi) \cap G_{\pi}\) has finite index in \(G_{\pi}\). If \(\chi\) is not trivial on \(G_{\pi}\) then the action of \(J\) on \(V_{\pi}\) is reducible. Indeed, \(\mathrm{ind}_{J}^{G_{\pi}}(1)\) contains subrepresentations \(1\) and \(\chi|_{G_{\pi}}\), and by Frobenius reciprocity \(\mathrm{End}_J(V_{\pi}|_J)\) is equal to \(\mathrm{Hom}_{G_{\pi}}(V_{\pi}, \mathrm{ind}_J^{G_{\pi}}(V_{\pi}|_J)) = \mathrm{Hom}_{G_{\pi}}(V_{\pi}, V_{\pi} \otimes \mathrm{ind}_J^{G_{\pi}}(1))\).

Hence \(\dim(\mathrm{End}_{J}(V_{\pi}|_J)) \ge 2\) and \(V_{\pi}|_J\) is reducible. By the hypothesis of multiplicity 1, \(V_{\pi}|_H\) is irreducible, hence \(V_{\pi}|_J\) is irreducible as \(H \subset J\). So \(\chi\) is trivial on \(G_{\pi}\).

The group \(G_{\pi}\) is a subgroup \(I\) of \(G\) containing \(H\) with \(\mathrm{len}(V|_I) = \mathrm{len}(V|_H)\). If \(I\) has this property, the restriction to \(H\) of any irreducible component on \(V|_I\) is irreducible, hence \(I\) is contained in \(G_{\pi}\).

\item[(d)] follows from (c).
\end{enumerate}
\end{enumerate}
\end{Proof}

% \begin{titledBox}{Restriction and Multiplicity}{resToMult}
%   Assume that \(V|_H\) has multiplicity 1.
%   The \(G\)-stabilizer of the isomorphism class of the irreducible components of \(V|_H\) is \(G_{\pi}\).
%   Let \(I\) be a subgroup of \(G\) such that \(H \subseteq I \subseteq G\).
%
%   Since \(G/H\) is abelian, \(G_\pi\) is normal. The length of the restriction to \(I\) is determined by the index of the composite group \(IG_\pi\):
%   \[ \mathrm{len}(V|_I) = [G : I G_\pi] \]
%
%   \textbf{Deduction for Prime Index:}
%   Suppose \([G:I] = p\) is a prime number (so \(I\) is maximal).
%   \begin{itemize}
%       \item If \(V|_I\) is \textbf{reducible}, then \(G_\pi \subseteq I\).
%       \item If \(V|_I\) is \textbf{irreducible}, then \(G_\pi \not\subseteq I\) (and \(IG_\pi = G\)).
%   \end{itemize}
% \end{titledBox}

Finally, with these results specifically for the case where the multiplicity is $1$, we need a way to detect
it. To that end, recall that the functor $\mathrm{Hom}_U(-, \theta)$ is always left-exact for a representation
$\theta$ of $H$ (this is even an exercise in Dummit and Foote).

\begin{lem}{Multiplicity One Criterion}{lem3}
  Let $V$ be an irreducible $K$-representation of $G$, $H$ a closed subgroup, and assume $V|_H$ has finite length.
  Let $\theta$ be a smooth $K$-representation of a closed subgroup $U \subseteq H$.

  \qn
  If the functor $\mathrm{Hom}_U(-, \theta)$ is exact on the category of smooth $K$-representations of $H$ and
  \[ \dim_K \mathrm{Hom}_U(V|_H, \theta) = 1, \]
  then $V|_H$ has multiplicity 1 (i.e., every irreducible constituent appears exactly once).
\end{lem}

\begin{Proof}
  Since $V|_H$ has finite length, let $0 = V_0 \subset V_1 \subset \dots \subset V_n = V|_H$ be a composition series where each quotient $W_i = V_i/V_{i-1}$ is irreducible.
  Let $\pi$ denote an isomorphism class of irreducible representations of $H$. Let $m_V(\pi)$ be the multiplicity of $\pi$ in this series (the number of indices $i$ such that $W_i \cong \pi$).

  By the exactness of the functor $\mathrm{Hom}_U(-, \theta)$, the dimension is additive:
  \[
    \dim \mathrm{Hom}_U(V|_H, \theta) = \sum_{i=1}^n \dim \mathrm{Hom}_U(W_i, \theta)
  \]
  Grouping by isomorphism class and using the hypothesis:
  \[
    1 = \sum_{\pi} m_V(\pi) \dim \mathrm{Hom}_U(\pi, \theta)
  \]
  Since $m_V(\pi)$ and $\dim \mathrm{Hom}_U(\pi, \theta)$ are non-negative integers, there must be exactly one isomorphism class $\pi_0$ such that:
  \[ m_V(\pi_0) = 1 \quad \text{and} \quad \dim \mathrm{Hom}_U(\pi_0, \theta) = 1 \]
  and for all other $\pi \neq \pi_0$, either the multiplicity is 0 or the Hom-dimension is 0.

  However, by Clifford's Theorem (which applies because $V$ is irreducible on $G$), the multiplicity $m_V(\pi)$ is constant for all irreducible constituents appearing in $V|_H$. Since we found one constituent ($\pi_0$) with multiplicity 1, all constituents must have multiplicity 1.
\end{Proof}

\section{p-adic Reductive Group}

We now focus in on when $G$ is a $p$-adic reductive group (like $\GL_n(F)$) and ensure that it satisfies the result we want.
% Insure to incorporate this: \url{https://en.wikipedia.org/wiki/Reductive_group}
Let $G$ be a $p$-adic (algebraic) reductive group, that is the group of $F$-rational points of a reductive
connected $F$-group $\underline{G}$, so $G= \underline{G}(F)$. We first ensure that the inclusion is
``well-defined'', as in it is closed and hence we have that $G/f(H)$ will be Hausdorff and $f(H)$ is locally
profinite. In particular, by our earlier work we will choose it so that it is compact abelian.

\begin{lem}{image is closed}{lem:3.1}
    Let $f : \underline{H} \to \underline{G}$ be an $F$-morphism of reductive connected $F$-groups. Then the subgroup $f(H)$ of $G$ is closed.
\end{lem}

\begin{Proof}
  see~\cite{henniartRepresentationsSL2F2025}. Note this is \emph{not} obvious: in general it is constructable
  (i.e. finite union of locally closed sets). The proof first shows it is locally closed, and then shows it
  has to be closed.
\end{Proof}


\begin{thm}{``Pushforward'' of representations}{3.2}
  Let \(f: \underline{H} \to \underline{G}\) be an \(F\)-morphism of reductive connected \(F\)-groups such
  that \(f(H)\) is a normal subgroup of \(G\) of compact quotient \(G/f(H)\). Then:

  \begin{enumerate}
    \item  The restriction to \(f(H)\) of any irreducible admissible \(K\)-representation of \(G\) is semisimple of finite length.
    \item Any irreducible admissible \(K\)-representation of \(f(H)\) is contained in some irreducible
      admissible \(K\)-representation of \(G\) restricted to \(f(H)\), and extends to an irreducible
      admissible representation of some open subgroup of \(G\) of finite index.
  \end{enumerate}
\end{thm}

\begin{Proof}
  We go over why the assumptions of~\cref{pr:admissibleExt} are true. Recall:
  \begin{enumerate}[label=(2-\arabic*)]
      \item Any finitely generated admissible $R$-representation of $G$ has finite length.
      \item Any finitely generated smooth $R$-representation of $H$ is Noetherian.
  \end{enumerate}
  In general, it is \emph{not} true that a locally profinite group satisfies (2-1). However, a reductive
  connected p-adic group does. This is too much to prove, however we have seen in~\cite{NatePadicRep} for
  $\GL_n(F)$ that supercuspidal have length $1$ and principal series have length at most $3$, and hence
  satisfy the condition.

  Next, $f(H)$ satisfies (2-2), but this is now harder to show and is a result of \emph{Bernstein's Theorem}.
  The high-level intuition is that the Hecke algebras $\mathcal{H}(G, K_0)$ are Noetherian and $\mathcal{H}(G)$ is built from
  these (using the Bernstein Blocks we talked about in class). Needless to say, we are taking this for
  granted.
\end{Proof}

\begin{prop}{}{prop:3.3}
    Let $f : H \to G$ be a surjective central\footnote{The kernel of the map is contained in the center of $H$} $F$-morphism of connected reductive $F$-groups. Then, the subgroup $f(H)$ of $G$ is normal of abelian compact quotient $G/f(H)$.
\end{prop}

\begin{Proof}
  Complicated, see~\cite{henniartRepresentationsSL2F2025}.
    % There is an $F$-morphism $\kappa:G\times G \to H$ such that $\kappa(f(x), f(y))=xhx^{-1}y^{-1}$ for all $x, y \in H$ \textcolor{red}{[Borel and Tits 1972, définition 2.2]}. So for all $u, v \in G$ we have $uvu^{-1}v^{-1} = f \circ \kappa(u,v) \in f(H)$. The subgroup $f(H)$ of $H$ is closed (Lemma 3.1) and normal with abelian quotient $G/f(H)$ \textcolor{red}{[loc. cit., proposition 2.7]}.

    % The compactness of $G/H$ is stated in \textcolor{red}{[Silberger 1979]} without proof and in \textcolor{red}{[Labesse and Schwermer 2019, Proposition A.2.1]} with indications for the proof. The idea is to reduce to a connected reductive $F$-anisotropic modulo the centre $F$-group.
    %
    % Let $S$ be a maximal $F$-split subtorus of $G$, and $B$ a parabolic $F$-subgroup of $G$ containing $S$. The $G$-centralizer $M$ of $S$ is compact modulo its centre and is a Levi component of $B$. Let $U$ be the unipotent radical of $B$. By \textcolor{red}{[Borel 1991, Theorem 22.6]}, the inverse image $S'$ of $S$ in $H$ is a maximal $F$-split torus in $H$, and the inverse image $B'$ of $B$ is a parabolic $F$-subgroup of $H$. Put $M'$ for the $H$-centralizer of $S'$ and $U'$ for the unipotent radical of $B'$. From \textcolor{red}{[loc. cit.]}, $f$ induces a surjective central $F$-morphism $M' \to M$ and an $F$-isomorphism $U' \to U$. On the other hand, we have the Iwasawa decomposition $G=KB$ for an open compact subgroup $K$ of $G$. The product map $K \times B \to G$ gives a surjective map $K \times B/f(B') \to G/f(H)$. We have $B/f(B') = M/f(M')$, so we just need to prove the compactness of $M/f(M')$.
    %
    % Let $X^*(S)$ denote the group of algebraic characters of $S$, and $\mathcal{S}(p_F)$ denote $\mathrm{Hom}(X^*(S), p_F^{\mathbb{Z}})$. The subgroup $\mathcal{S}(p_F)$ of $S$ is free abelian of finite rank with a compact quotient $S/\mathcal{S}(p_F)$. On the other hand, $f$ induces a surjective $F$-morphism $S' \to S$ sending $S'(p_F)$ onto a sublattice of $\mathcal{S}(p_F)$. Hence $S/f(S')$ is finite. So $M/f(S')$ is compact since $M/S$ is compact, a fortiori $M/f(M')$ is compact.
\end{Proof}

  The condition that the morphism $f$ is \emph{central} is necessary. Without it, the image may fail to be normal or co-compact.

\begin{example}{Necessity of Central}{necessityofCentral}
  Assume $\mathrm{char}(F) = 2$. Let $\underline{G} = \underline{\mathrm{GL}}_2$ and $\underline{G}' = \underline{\mathrm{SL}}_2$.
  Consider the Frobenius map $\varphi(x) = x^2$ (applied entry-wise to matrices).
  Define the morphism $f: \underline{G} \to \underline{G}'$ by:
  \[ f(g) = \varphi(g) \cdot (\det g)^{-1} \]
  (Note: In char 2, $\det(\varphi(g)) = (\det g)^2$, so $\det(f(g)) = (\det g)^2 (\det g)^{-2} = 1$, so the map lands in $\mathrm{SL}_2$).

  This map is surjective as algebraic groups, but it is \emph{not central}.
  Let $U$ be the subgroup of unipotent upper-triangular matrices in $\mathrm{SL}_2(F)$ (isomorphic to the additive group of $F$).
  The image of $U$ under $f$ is $\varphi(U)$, which corresponds to the subgroup of squares $F^2 \subset F$.
  \begin{enumerate}
      \item \emph{Not Normal:} The image $f(G(F))$ involves matrices with entries in $F^2$. Conjugating a matrix with square entries by a matrix with non-square entries generally results in non-square entries.
      \item \emph{Not Cocompact:} The quotient group contains the term $U / \varphi(U)$, which is homeomorphic to the additive quotient $F / F^2$. For a local field of characteristic 2 (like $\mathbb{F}_{2^k}((t))$), the space $F/F^2$ is \emph{not compact} (it is essentially an infinite-dimensional vector space or a copy of the field itself).
  \end{enumerate}
\end{example}


\begin{cor}{Descent via Derived Groups}{3.5}
Assume \(K\) is algebraically closed. Let \(f: \underline{H} \to \underline{G}\) be an \(F\)-morphism of connected reductive \(F\)-groups.
Suppose that \(f\) induces a \emph{central isogeny} on the derived groups:
\[ f^{\mathrm{der}}: \underline{H}^{\mathrm{der}} \to \underline{G}^{\mathrm{der}} \]
Then the subgroup \(f(H) \subset G\) satisfies the conclusions of \cref{th:3.2} (i.e., restrictions of admissible representations are semisimple of finite length, and admissible extensions exist).
\end{cor}

\begin{Proof}
    \textbf{Step 1: Covering the Group.}
    By definition of a central isogeny, the map \(f^{\mathrm{der}}\) is surjective with a finite kernel contained in the center of \(\underline{H}^{\mathrm{der}}\) \textcolor{red}{[Springer 1998, §12.2.6]}.
    Let \(\underline{Z}\) be the connected center of \(\underline{G}\). Structural theory of reductive groups tells us that \(\underline{G}\) is generated by its center and its derived group. Specifically, the multiplication map
    \[ m: \underline{Z} \times \underline{G}^{\mathrm{der}} \to \underline{G} \]
    is surjective \textcolor{red}{[Springer 1998, Corollary 8.1.6]}.

    \textbf{Step 2: Constructing a Central Morphism.}
    Consider the map \(\Psi: \underline{Z} \times \underline{H} \to \underline{G}\) defined by \(\Psi(z, h) = z \cdot f(h)\).
    Since \(f\) covers \(\underline{G}^{\mathrm{der}}\) and we explicitly included \(\underline{Z}\), \(\Psi\) is surjective.
    Its kernel lies in the product of the center of \(\underline{Z}\) (which is \(\underline{Z}\) itself) and the kernel of \(f\), which is central. Thus, \(\Psi\) is a surjective central morphism.

    \textbf{Step 3: Applying Proposition 3.3.}
    By \cref{pr:prop:3.3}, the image \(\Psi(Z \times H) = Z \cdot f(H)\) is a normal subgroup of \(G\) with a compact abelian quotient.
    Consequently, by \cref{th:3.2}, the restriction of any irreducible representation \(V\) of \(G\) to the subgroup \(Z \cdot f(H)\) has finite length and is semisimple.

    \textbf{Step 4: Removing the Center (Schur's Lemma).}
    Since \(V\) is irreducible and \(K\) is algebraically closed, by Schur's Lemma, the center \(Z\) acts on \(V\) via a scalar character \(\omega: Z \to K^\times\).
    Therefore, the submodule structure of \(V\) restricted to \(Z \cdot f(H)\) is identical to the submodule structure restricted to just \(f(H)\) (since the action of \(Z\) preserves all subspaces).
    Thus, \(V|_{f(H)}\) is also semisimple of finite length.
\end{Proof}

\section{Restriction of $GL_2(F)$ to $SL_2(F)$}


Let $G = \GL_2(F)$, $G' =\SL_2(F)$, $Z$ the center, $T$ the torus, $B$ the Borel subgroup, $U$ the unipotent
subgroup. If $X \sse G$ is a subset, $X' = X\cap G'$ and $x = (x_{ij})$ represents a matrix in $G$. We first
set up for the Clifford theory.

Take $H= ZG'$. It is a closed normal compact group where $G/H$ is isomorphic to $F^*/(F^*)^2$. It is
normal as $Z$ is obviously normal and $G' = \ker\det$ where $\mathrm{im}\det = F^*$, and $gZG'g^{-1} = ZG'$. It is
closed since $H= \det^{-1}\left((F^*)^2\right)$. It is cocompact since:
\[
  \frac{F^*}{\left(F^*\right)^2} \cong \frac{\Z}{2\Z} \times \frac{\CO_F^*}{(\CO_F^*)^2}
\]
and $\CO_F^*$ is compact and so is the quotient by a closed subgroup The $\mathbb{F}_2$-dimension of $G/H$ is that of
$F^*/(F^*)^2$ which is:
\begin{equation}\label{eq:4-1}
\dim_{\mathbb{F}_2} F^* / (F^*)^2 =
\begin{cases}
2 + e & \text{if } \text{char}_F \neq 2, \\
\infty & \text{if } \text{char}_F = 2,
\end{cases}
\quad \text{where } 2\CO_F = \FP_F^e.
\end{equation}
Hence if $\operatorname{char}_F \ne 2$, $ZG'$ is open.


Next, we set up to show that the multiplicity of the restriction is $1$. Fix a separable closure $F^{\mathrm{sc}}$ of $F$ and will consider only extensions of $F$ contained in
$F^{\mathrm{sc}}$. Write $W_F$ for the Weil group and $\mathrm{Gal}_F$ for the Galois group of
$F^{\mathrm{sc}}/F$. Let $\Psi$ be a smooth additive $K$-character of $F$, trivial on $\CO_F$ but not on
$\mathfrak{p}_F^{-1}$.

\begin{defn}{Whittaker Space}{whittakerSp}
  The smooth $K$-characters of $U$ have the form
  \[
    \theta_Y(u) = \Psi \circ \tr(Y(u-1)) = \Psi(Y_{2,1}u_{1,2}), \qquad u\in U
  \]
  for a lower triangular nilpotent matrix $Y \in M_2(F)$. The case $Y=0$ gives the trivial character of $U$,
  the cases with $Y\ne 0$ give the \emph{nondegenerate characters} of $U$.

  \qn
  When $Y_{2,1} = 1$, we denote $\theta_Y = \theta$.
\end{defn}


The normalizer of $U$ in $G$ is $TU$. By conjugation, $U$ acts trivially on $U$ and its characters. A diagonal matrix $t = \mathrm{diag}(t_1, t_2)$ acts on $u \in U$ by $(tut^{-1})_{1,2} = (t_1/t_2)u_{1,2}$. Also, $t$ acts on the matrix $Y$ defining the character by $(tYt^{-1})_{2,1} = (t_2/t_1)Y_{2,1}$.

It follows that $T$ acts transitively on the nondegenerate characters of $U$, the quotient $T/Z$ acting simply transitively. By the same formulas, two nontrivial characters $\theta_Y$ and $\theta_{Y'}$ of $U$ are conjugate in $G'$ if and only if they are conjugate by an element of $T'$ (the diagonal torus in $\SL_2$) and if and only if $Y_{1,2}$ and $Y'_{1,2}$ differ by a square in $F^*$.

The $T$-normalizer of $\theta_Y$ is equal to $Z$ if $Y \neq 0$ and to $T$ if $Y=0$.
The $\theta_Y$-coinvariant functor $\tau \mapsto W_Y(\tau)$ from the smooth $K$-representations $\tau$ of $U$
to the smooth $K$-representations of the $T$-normalizer of $\theta_Y$ is exact. This functor is essentially the dual of $\operatorname{Hom}_U(\tau, \theta)$.
\begin{itemize}
    \item A smooth $K$-representation $\tau$ of $U$ is called \emph{degenerate} when $W_Y(\tau) = 0$ for all $Y \neq 0$, and \emph{nondegenerate} otherwise.
    \item A smooth representation of $G$ or $G'$ is called degenerate (or nondegenerate) if its restriction to $U$ is as well.
\end{itemize}

The finite-dimensional irreducible smooth $K$-representations of $G$ are of the form $\chi \circ \det$ for a smooth character $\chi$ of $F^*$ and are degenerate. If $\Pi$ is an infinite-dimensional irreducible smooth $K$-representation of $G$, then:
\[ \dim W_Y(\Pi) = 1 \]
for all $Y \neq 0$ by the uniqueness of Whittaker models \textcolor{red}{[Vignéras 1996, chap III, 5.10]} .

This uniqueness is structurally decisive. By applying the \emph{Multiplicity One Criterion} (\cref{lm:lem3}), we conclude that the restriction $\Pi|_{G'}$ is \emph{multiplicity free} (i.e., it
decomposes into a direct sum of \emph{distinct} irreducible representations). Moreover, by Clifford's Theorem,
these components form a single orbit under the conjugation action of $G$. This natural grouping of ``conjugate
siblings'' forms the basis of the classification for $G'$.

\index{L-packet}
\begin{defn}{L-packet}{l-packet}
Let \(\Pi\) be an irreducible smooth \(K\)-representation of \(G\). The \emph{L-packet} associated to \(\Pi\), denoted \(L(\Pi)\), is the set of isomorphism classes of the irreducible components of the restriction \(\Pi|_{G'}\).
\end{defn}

Since every irreducible representation of \(G'\) can be lifted to \(G\) (by~\cref{pr:admissibleExt}), we can classify the irreducible smooth \(K\)-representations of \(G'\) by analyzing these L-packets.

\begin{prop}{Properties of L-packets: I}{propofL-packetsI}
  \begin{enumerate}
    \item If \(\Pi\) is an irreducible smooth \(K\)-representation of \(G\), the restriction \(\Pi|_{G'}\)
    is semisimple of finite length and has multiplicity 1.
    \item Any irreducible smooth \(K\)-representation of \(G'\) belongs to a unique \(L\)-packet.
    \item For two irreducible smooth \(K\)-representations \(\Pi_1, \Pi_2\) of \(G\):
    \[
      L(\Pi_1) = L(\Pi_2) \iff \Pi_1 \cong (\chi \circ \det) \otimes \Pi_2
    \]
    for some smooth character \(\chi\) of \(F^*\).
  \end{enumerate}
\end{prop}

\begin{Proof}
  \begin{enumerate}
    \item \textbf{Semisimplicity:} We established in the previous section that \(H = ZG'\) is a normal subgroup with a compact abelian quotient. By \cref{th:3.2} (Descent Properties), the restriction \(\Pi|_H\) is semisimple of finite length. Since \(Z\) acts by scalars (as \(K\) is algebraically closed), the restriction to \(G'\) retains the same reducible structure.

    \textbf{Multiplicity 1:} We apply the Multiplicity One Criterion (\cref{lm:lem3}). Let \(U\) be the unipotent subgroup and \(\theta\) a non-degenerate character. By the uniqueness of Whittaker models for \(\mathrm{GL}_2(F)\), for any infinite-dimensional \(\Pi\), we have \(\dim \mathrm{Hom}_U(\Pi, \theta) = 1\). Thus, \(\Pi|_{G'}\) is multiplicity free. (For the finite-dimensional case, \(\Pi\) is 1-dimensional and trivial on \(G'\), so the result holds trivially).

    \item \textbf{Existence:} By the Extension Property (\cref{pr:admissibleExt}), any irreducible representation \(\pi\) of \(G'\) is contained in the restriction of some representation of \(G\). Thus, \(\pi\) belongs to at least one L-packet.
    \textbf{Uniqueness:} By Clifford's Theorem, the irreducible components of a restriction \(\Pi|_{G'}\) form a single orbit under the conjugation action of \(G\). By Definition \ref{df:l-packet}, the L-packet \(L(\Pi)\) is exactly this orbit. Since distinct orbits are disjoint, \(\pi\) cannot belong to two different packets.

    \item
    \begin{itemize}
        \item[($\Leftarrow$)] Assume \(\Pi_1 \cong \Pi_2 \otimes (\chi \circ \det)\). Restricting to \(G'\):
        \[ \Pi_1|_{G'} \cong (\Pi_2|_{G'}) \otimes (\chi \circ \det)|_{G'} \]
        Since \(\det(G') = \{1\}\), the twisting character is trivial on \(G'\). Thus \(\Pi_1|_{G'} \cong \Pi_2|_{G'}\), which implies their components (and thus their L-packets) are identical.

        \item[($\Rightarrow$)] Assume \(L(\Pi_1) = L(\Pi_2)\). This implies that \(\Pi_1|_{G'}\) and \(\Pi_2|_{G'}\) share an irreducible component \(\pi\).
        By the logic of \cref{lm:lem2} (Clifford Theory for Abelian Quotients). Although \(G/G' \cong F^*\) is not compact, we can pass through the intermediate compact subgroup \(H = ZG'\).
        Since \(\Pi_1\) and \(\Pi_2\) share a component on \(G'\), they must share a component on \(H\) (possibly up to a twist of the central character).
        Since \(G/H\) is abelian, \cref{lm:lem2}(2b) implies that \(\Pi_1\) is a twist of \(\Pi_2\) by some character \(\psi\) of \(G\) trivial on \(H\).
        Since \(\det\) gives the isomorphism \(G/G' \cong F^*\), any such twist is of the form \(\chi \circ \det\).
    \end{itemize}
  \end{enumerate}
\end{Proof}

\begin{titledBox}{remark}{rmk}
  The trivial character of \(G'\) is the unique finite-dimensional irreducible smooth \(K\)-representation of \(G'\). It is degenerate. It forms the L-packet \(L(\mathbf{1}) = L(\chi \circ \det)\) for any smooth character \(\chi\) of \(F^*\).
\end{titledBox}

Thus, the irreducible components of \(\Pi|_{G'}\) are \(G\)-conjugate (and indeed \(B\)-conjugate, since the Iwasawa decomposition \(G = BG'\) implies the Borel subgroup covers the quotient).

Let us now determine the size of an \(L\)-packet. As established in \cref{lm:lem3}, since \(G/G'\) is abelian,
the stabilizer \(G_{\pi}\) of a component \(\pi \in L(\Pi)\) is a normal subgroup and does not depend on the
choice of \(\pi\). Let us denote this as \(G_{\Pi}\). Since \(G_{\Pi}\) contains the center \(Z\), its image
under the determinant \(\det(G_{\Pi})\) must contain \(\det(Z) = (F^*)^2\). We have the following isomorphism
of quotients:
\[
  G/G_\Pi \cong \frac{G/G'}{G_\Pi/G'} \xrightarrow{\sim} \frac{F^*}{\det(G_\Pi)}
\]
This isomorphism allows us to count the size of the packet using the arithmetic of the field \(F\).

\begin{prop}{Properties of L-packets: II}{propofL-packetsII}
  Let \(\Pi\) be an irreducible smooth \(K\)-representation of \(G\) and \(\pi \in L(\Pi)\) be an irreducible component. Let \(G_\Pi\) denote the \(G\)-stabilizer of \(\pi\).
  \begin{enumerate}
    \item \textbf{Maximality:} \(G_\Pi\) is the largest subgroup \(I\) of \(G\) (containing \(G'\)) such that \(\mathrm{length}(\Pi|_I) = \mathrm{length}(\Pi|_{G'})\). (In fact, \(\mathrm{length}(\Pi|_{G_\Pi}) = 1\)).
    \item \textbf{Induction:} \(\Pi \cong \operatorname{Ind}_{G_\Pi}^G \pi\). (Since multiplicity is 1, the isotypic component \(V_\pi\) is isomorphic to \(\pi\)).
    \item \textbf{Structure:} The L-packet \(L(\Pi)\) is a principal homogeneous space (a \emph{torsor}) for the quotient group \(G/G_\Pi\).
    \item \textbf{Cardinality:} \(|L(\Pi)|\) is a power of 2. Specifically, \(|L(\Pi)|\) divides the index \([F^* : (F^*)^2]\).
    \[
      |L(\Pi)| \text{ divides } \begin{cases}
        2^{2+e} & \text{if } \mathrm{char}_F \neq 2, \\
        \infty & \text{if } \text{char}_F = 2.
      \end{cases}
    \]
  \end{enumerate}
\end{prop}


\begin{Proof}
  \begin{enumerate}
    \item By \cref{pr:propofL-packetsI}, the restriction \(\Pi|_{G'}\) has multiplicity 1. By \cref{lm:lem3}, specifically part (2c), \(G_\Pi\) is the unique maximal subgroup \(I\) such that \(\mathrm{length}(\Pi|_I) = \mathrm{length}(\Pi|_{G'})\).
    (Note: While the lemma refers to \(H = ZG'\), the scalar action of \(Z\) ensures that lengths on \(H\) and \(G'\) are identical).

    \item Since multiplicity is 1, the isotypic component \(V_\pi\) coincides with the irreducible
      representation \(\pi\). The isomorphism \(\Pi \cong \operatorname{Ind}_{G_\Pi}^G \pi\) is then the
      standard Clifford reduction formula established in the setup to \cref{lm:lem3}.

    \item Clifford theory establishes that \(G\) acts transitively on \(L(\Pi)\) with stabilizer \(G_\Pi\). Since \(G/G'\) is abelian, \(G_\Pi\) is normal. Thus, the quotient group \(G/G_\Pi\) acts simply transitively on the set \(L(\Pi)\), making it a \(G/G_\Pi\)-torsor.

    \item The cardinality is the index: \(|L(\Pi)| = [G : G_\Pi]\).
    The determinant induces an isomorphism \(G/G_\Pi \cong F^*/\det(G_\Pi)\).
    Since \(G_\Pi\) contains the center \(Z\), its determinant image contains \(\det(Z) = (F^*)^2\).
    Therefore, the group \(F^*/\det(G_\Pi)\) is a quotient of the elementary abelian 2-group \(F^*/(F^*)^2\).
    By Lagrange's Theorem, the size of this quotient must be a power of 2 dividing \([F^* : (F^*)^2]\).
  \end{enumerate}
\end{Proof}



\begin{thm}{Size of L-packets (Odd Residual Characteristic)}{sizeofL-packets}
  If the residual characteristic \(p\) is odd, then the cardinality of an \(L\)-packet is either \(1\), \(2\), or \(4\).
\end{thm}

\begin{Proof}
  According to \cref{pr:propofL-packetsII}(4), \(|L(\Pi)|\) divides the order of the group \(F^*/(F^*)^2\).
  Recall Equation \eqref{eq:4-1}, which states that \(\dim_{\mathbb{F}_2} F^*/(F^*)^2 = 2 + e\), where \(e = v_F(2)\).
  Since \(p\) is odd, \(2\) is a unit in the ring of integers \(\CO_F\), implying \(e = 0\).
  Therefore, the group has dimension 2 over \(\mathbb{F}_2\), so \(|F^*/(F^*)^2| = 2^2 = 4\).
  The divisors of 4 are 1, 2, and 4.
\end{Proof}

When $p =2$, this is more difficult (If $F = \mathbb{Q}_2$, the group size is $|F^*/(F^*)^2| = 8$). However, the
packets are \emph{still} of size $1,2,4$ as long as $\operatorname{char}_K \ne 2$. We shall not prove this, but it was proven
for:
\begin{itemize}
  \item $K = \mathbb{C}$ and $\operatorname{char}_F = 0$ (see \textcolor{red}{[Labesse and Langlands 1979; Shelstad 1979]}
  \item $\operatorname{char}_F = 2$ and $\operatorname{char}_K = 2$ \textcolor{red}{[Cui et al. 2024, Proposition 6.6]}
\end{itemize}

The isomorphism \(G/G_\Pi \cong F^*/\det G_\Pi\) invites an interpretation via Local Class Field Theory. Recall that open subgroups of \(F^*\) of finite index correspond to norm groups of abelian extensions.
\begin{itemize}
    \item Any open subgroup of index 2 is the norm group \(N_{E/F}(E^*)\) for a unique quadratic separable extension \(E/F\).
    \item Any open subgroup of index 4 containing \((F^*)^2\) is the norm group \(N_{K/F}(K^*)\) for a unique biquadratic separable extension \(K/F\).
\end{itemize}

\begin{defn}{Associated Field Extension}{fieldExt}
  Let \(\Pi\) be an irreducible smooth \(K\)-representation of \(G\). We denote by \(E_\Pi\) the separable extension of \(F\) corresponding to the stabilizer via Class Field Theory, i.e., the extension such that:
  \[ N_{E_\Pi/F}(E_\Pi^*) = \det G_\Pi \]
\end{defn}

We can now relate the L-packet size to the group of self-twists of \(\Pi\).

\begin{defn}{Group of Self-Twists}{selfTwists}
  We denote by \(X_\Pi\) the group of characters \(\chi \circ \det\) of \(G\) (where \(\chi\) is a smooth character of \(F^*\)) such that:
  \[ \Pi \otimes (\chi \circ \det) \cong \Pi \]
  If \(\mathrm{char}(K) \neq 2\), the non-trivial characters in \(X_\Pi\) correspond to quadratic characters \(\eta_E\) associated with quadratic extensions \(E/F\).
\end{defn}

\begin{prop}{Twists and Packet Size}{twistsPacketSize}
  \begin{enumerate}
      \item \(X_\Pi\) is isomorphic to the group of smooth characters of \(G\) trivial on \(G_\Pi\).
      \item The cardinality of the L-packet equals the number of self-twists: \(|L(\Pi)| = |X_\Pi|\).
  \end{enumerate}
\end{prop}

\begin{Proof}
  \begin{enumerate}
      \item This is a direct application of \cref{lm:lem3}(1). Since \(G/H\) is abelian, the self-twists of \(\Pi\) are exactly the characters trivial on the stabilizer \(G_\Pi\).
      \item By Pontryagin duality  for the finite abelian group \(G/G_\Pi\)
        (see~\cite{chwojko-srawleyPontryaginDuality}), the number of characters trivial on \(G_\Pi\) equals
        the index \([G:G_\Pi]\). By \cref{pr:propofL-packetsII} we have \(|L(\Pi)| = [G:G_\Pi]\).
  \end{enumerate}
\end{Proof}

Finally, we address how the Whittaker model interacts with the L-packet. While all components are conjugate, they are distinct with respect to a fixed Whittaker character.

\begin{prop}{Uniqueness of Whittaker Component}{whittakerComponent}
  Fix a lower triangular matrix \(Y \neq 0\) and the associated character \(\theta_Y\) of \(U\).
  There exists a \emph{unique} component \(\pi \in L(\Pi)\) such that \(\pi\) admits a non-zero Whittaker model for \(\theta_Y\) (i.e., \(W_Y(\pi) \neq 0\)).
\end{prop}

\begin{Proof}
  Recall that \(\Pi|_{G'}\) decomposes as a direct sum of the components in the L-packet:
  \[ \Pi|_{G'} = \bigoplus_{\pi \in L(\Pi)} \pi \]
  The Whittaker functor \(\tau \mapsto W_Y(\tau)\) is exact and additive on direct sums. Therefore:
  \[ \dim_K W_Y(\Pi) = \sum_{\pi \in L(\Pi)} \dim_K W_Y(\pi) \]
  We know from the uniqueness of Whittaker models for \(\mathrm{GL}_2\) that \(\dim_K W_Y(\Pi) = 1\).
  Since dimensions are non-negative integers, exactly one term in the summation must be 1, and the rest must be 0.
\end{Proof}



\section{Representations}

We now briefly take a look at how we translate the result about parabolic induction and cuspidal
representations. Reminders of the definitions will be put, but it will be assumed the reader already knows
this (see~\cite{NatePadicRep} for my covering of this material). For this section, let \(\mathrm{Gr}_K^\infty(G)\) denote the Grothendieck group of finite
length smooth \(K\)-representations of \(G\), and let \([\tau]\) denote the image in
\(\mathrm{Gr}_K^\infty(G)\) of a representation \(\tau\). Similar notation applies to \(G'\).


\begin{defn}{Parabolic Induction}{parabolicInd}
  Let \((\sigma, V)\) be a smooth \(K\)-representation of the torus \(T\). The \emph{smooth parabolic induction} \(\operatorname{ind}_B^G(\sigma)\) is the space of functions \(f: G \to V\) satisfying:
  \[
    f(tugk) = \sigma(t)f(g)
  \]
  for all \(t \in T, u \in U, g \in G\), and some open compact subgroup \(K_f \subset G\), under the action of right translation.


  The modulus character \(\delta: B \to \mathbb{R}_{>0}\) is defined by:
  \[
    \delta(\mathrm{diag}(a, d)) = q^{-\mathrm{val}(a/d)} \qquad (a, d \in F^*)
  \]
  To define the square root of \(\delta\), fix a square root \(q^{1/2}\) of \(q\) in \(K^*\):
  \[
    \nu(\mathrm{diag}(a, d)) = (q^{1/2})^{-\mathrm{val}(a/d)}
  \]
  Then the \emph{normalized parabolic induction} is defined as \(i_B^G(\sigma) = \operatorname{ind}_B^G(\sigma \nu)\).
\end{defn}

This construction behaves well with respect to twisting. For a smooth character \(\chi \circ \det\) of \(G\):
\[ i_B^G(\sigma) \otimes (\chi \circ \det) \simeq i_B^G(\sigma \otimes (\chi \circ \det)) \]
Analogous definitions apply to \(G'\). Since \(G = BG'\) and \(G/B\) is compact, the restriction map \(f \mapsto f|_{G'}\) induces isomorphisms:
\[
  (\operatorname{ind}_B^G(\sigma))|_{G'} \xrightarrow{\sim} \operatorname{ind}_{B'}^{G'}(\sigma|_{T'}), \qquad (i_B^G(\sigma))|_{G'} \xrightarrow{\sim} i_{B'}^{G'}(\sigma|_{T'})
\]
Hence, the parabolic case is rather simple.
\begin{defn}{Cuspidal Representations}{cuspidalDef}
  When \(\chi\) is a smooth \(K\)-character of \(T\), the induction \(\operatorname{ind}_B^G(\chi)\) is called a \emph{principal series} representation.
  \begin{itemize}
      \item An irreducible smooth \(K\)-representation of \(G\) is called \emph{supercuspidal} if it is not a subquotient of any principal series.
      \item It is called \emph{cuspidal} if its \(U\)-coinvariants are zero (\((-)_U = 0\)).
  \end{itemize}
  A supercuspidal representation is always cuspidal. The converse holds provided \(q+1 \neq 0\) in \(K\). Both classes of representations are infinite-dimensional.
\end{defn}

The property of cuspidality descends to the L-packet.

\begin{prop}{Cuspidality of L-packets}{cuspidalPacket}
  Let \(\Pi\) be an irreducible smooth \(K\)-representation of \(G\) and \(\pi \in L(\Pi)\) a component of its restriction.
  \begin{enumerate}
      \item \(\Pi\) is cuspidal if and only if \(\pi\) is cuspidal.
      \item \(\Pi\) is supercuspidal if and only if \(\pi\) is supercuspidal.
  \end{enumerate}
  Consequently, an L-packet \(L(\Pi)\) is called a \emph{cuspidal packet} if \(\Pi\) is cuspidal.
\end{prop}

\begin{Proof}
  \(L(\Pi)\) constitutes the \(B\)-orbit of \(\pi\) (since \(G=BG'\)). The \(U\)-coinvariant functor is exact and commutes with restriction to \(G'\). Thus, the vanishing of coinvariants for \(\Pi\) implies the vanishing for its components, and vice versa.
\end{Proof}

\subsubsection*{Construction via Types}

The rest would now be going into the details of the construction which is too much for this paper, but here is
a glimpse. Let \(\Pi\) be a cuspidal \(K\)-representation of \(G\). It is constructed via compact induction from an \emph{extended maximal simple type} \((J, \Lambda)\):
\[
  \Pi \cong \operatorname{ind}_J^G(\Lambda)
\]
References: \textcolor{red}{[Bushnell and Kutzko 1994; Bushnell and Henniart 2002]} for \(K=\mathbb{C}\), and \textcolor{red}{[Vignéras 1996, III.3.4]} for general \(K\).

The group \(J\) contains the center \(Z\) and a unique maximal open compact subgroup \(J^0\). Let \(J^1\) be
the pro-\(p\) radical of \(J^0\). The restriction \(\Lambda|_{J^0}\) is of the form \(\kappa \otimes
\tilde{\sigma}\), where \(\kappa\) is a representation of \(J^0\) and \(\tilde{\sigma}\) is the inflation of
an irreducible representation \(\sigma\) of \(J^0/J^1\).

The open normal subgroup \(JG'\) has index \(|F^*/\det J|\) in \(G\).

\begin{prop}{Mackey Decomposition and Packet Size}{mackeyPacket}
  The restriction of \(\Pi\) to the intermediate group \(JG'\) decomposes as:
  \[ \Pi|_{JG'} = \bigoplus_{g \in G/JG'} \operatorname{ind}_{J_g'}^{JG'} \lambda^g \]
  where \(J' = J \cap G'\) and \(g\) runs over coset representatives. This leads to the formula for the packet size:
  \[ |L(\Pi)| = [F^* : \det J] \cdot \mathrm{length}(\operatorname{ind}_{J'}^{G'} (\lambda|_{J'})) \]
\end{prop}

\begin{prop}{Quadratic Extension Criterion}{quadExtCrit}
  If the index \([F^* : \det J] = 2\), then \(\det G_\Pi \subset \det J\).
  In this case, \(\det G_\Pi = \det J\) if and only if \(G_\Pi = JG'\).
  This determines a unique quadratic separable extension \(E/F\) such that \(\det J = N_{E/F}(E^*)\).
\end{prop}

\subsubsection*{Next Steps}

With the structural properties of L-packets established (cardinality, relationship to field extensions, and type theory), the natural progression is to classify these packets explicitly.
\begin{itemize}
    \item \textbf{Level 0:} Cases where the type \((J, \Lambda)\) has trivial pro-\(p\) part (depth zero).
    \item \textbf{Positive Level:} Cases where the representation is deeper.
\end{itemize}
The subsequent sections should define the specific types \((J, \Lambda)\) for the primitive (cuspidal)
representations and calculate the corresponding groups \(X_\Pi\) and stabilizer groups \(G_\Pi\) to verify the
cardinality bounds derived in \cref{pr:twistsPacketSize}.




\newpage
\printbibliography






\end{document}

